% !PRINTED 119
\weberpage{119}
%<<CONT
can denote. Here too the law
\weq{6}{\pi_a \pi_{a'} = \pi_{aa'},}
holds, if $a'$ is likewise any element from $P$. There remains
the question whether different elements $a$ can give rise to the
same permutation $\pi$. If $\pi_a = \pi_{a'}$, then it follows from (6):
\[
\pi_{a'a^{-1}} = 1,
\]
and if therefore we denote by $a_0$ all the elements satisfying
the condition
\weq{7}{\pi_{a_0} = 1}
then $a' = a_0 a$, and it comes down to determining the totality
of the elements $a_0$. But if
\weq{8}{Rba_0 = Rb}
is to hold, then
\[
ba_0 = cb
\]
or
\[
a_0 = b^{-1}cb
\]
must hold, i.\,e. $a_0$ must belong to the group $b^{-1}Rb$, and this is
also sufficient for (8) to hold.

From this it follows that the elements $a_0$ satisfying the
condition (7) fill the whole group $R_0$, which is the intersection
of all the divisors of $P$ conjugate to $R$:
\[
R,\ b_1^{-1}Rb_1,\ b_2^{-1}Rb_2,\ \ldots,\ b_{j-1}^{-1}Rb_{j-1},
\]
and, as we have proved earlier, is a normal divisor
of $P$.

From this there results, for the case that $R_0 = 1$, the theorem:

2.\ \gesperrt{If a group $P$ has a divisor of index $j$
which has no element other than the unit
element in common with its conjugate
divisors, then $P$ is (simply) isomorphic with a
transitive permutation group of $j$ figures.}

That the permutation group is transitive is seen from (4),
where each of the elements $b_1, b_2, \ldots, b_{j-1}$ can be put for $a$.

Theorem 2. always takes place when $P$ is a \gesperrt{simple group}
and $R$ a proper divisor of $P$, because then $R_0$, as a normal divisor of $P$, must reduce to the unit group\footnotemark.

Now let $R$ and $Q$ denote any two divisors of the group $P$,
$R$ of index $j$, $Q$ of degree $q$.
\footnotetext{\gesperrt{Hölder}, Mathematische Annalen, Bd.\ 40, S.\ 57.}
